% GENERATED FILE. Edit canonical lessons, then run npm run generate:books.
% canonical-source-hash: fnv1a64:bc89652ffb9960a1
% canonical-lessons: IT-C330-epoca, IT-C330-medioevo, IT-C330-rovine, IT-C330-scavi, IT-C330-archeologo

\chapter{Sotto la terra}
\label{ch:it-sotto-la-terra}
\hlchaptermodality{330}

\begin{chapteropening}
\textbf{By the end of this chapter:} \emph{I can talk about the past under the ground: the age, the Middle Ages, ruins, digs and the archaeologist.}

Complete the last lesson of chapter 330: I can talk about the past under the ground: the age, the Middle Ages, ruins, digs and the archaeologist.
\end{chapteropening}

\section[l'epoca]{l'epoca — the age, the era, the period}

\label{lesson:IT-C330-epoca}



\begin{quote}
Before the new one: say the Italian for to surrender, then the Italian for to free.
\end{quote}

\subsection*{You'll want to know: l'epoca}
\textbf{l'epoca} — \textquotedblleft{}the age, the era, the period\textquotedblright{}.

\begin{quote}
In quell'epoca non c'era la luce elettrica.
\end{quote}

\begin{quote}
\emph{In those days there was no electric light.}
\end{quote}

\emph{All'epoca} — at the time.

\begin{cousinweb}
\emph{epoca} is a learned word usually traced to Greek \emph{epokh\'{\={e}}} \textquotedblleft{}a pause, a fixed point in time\textquotedblright{}, taken in through Medieval Latin (EtymDB has no Italian entry).
\end{cousinweb}

\begin{grammarlens}[title={What you've built}]
One more word for telling, explaining and describing at B1.
\end{grammarlens}

\subsection*{Your turn}
\begin{itemize}
\raggedright
  \item \emph{Say it:} \emph{l'epoca}
  \item \emph{Say it:} \emph{l'epoca}, once more
  \item \emph{Say it:} a sentence of your own with \emph{l'epoca}
  \item \emph{Recall:} say the Italian for to surrender, then the Italian for to free, then say \emph{l'epoca} again
\end{itemize}

\begin{tcolorbox}[breakable,colback=teal!4,colframe=teal!35!black,title={Before you move on}]
What does l'epoca mean? (\textquotedblleft{}The age, the era, the period\textquotedblright{}.) Say it once more.
\end{tcolorbox}
\section[il medioevo]{il medioevo — the Middle Ages}

\label{lesson:IT-C330-medioevo}



\begin{quote}
Before the new one: say the Italian for to free, then the Italian for the age.
\end{quote}

\subsection*{You'll want to know: il medioevo}
\textbf{il medioevo} — \textquotedblleft{}the Middle Ages\textquotedblright{}.

\begin{quote}
Il castello è del medioevo.
\end{quote}

\begin{quote}
\emph{The castle dates from the Middle Ages.}
\end{quote}

Often written \emph{Medioevo}; the adjective is \emph{medievale}.

\begin{cousinweb}
\emph{medioevo} joins \emph{medio} \textquotedblleft{}middle\textquotedblright{}, from Latin \emph{medius}, and \emph{evo} \textquotedblleft{}age\textquotedblright{}, from Latin \emph{aevum} (Lewis \& Short).
\end{cousinweb}

\begin{grammarlens}[title={What you've built}]
One more word for telling, explaining and describing at B1.
\end{grammarlens}

\subsection*{Your turn}
\begin{itemize}
\raggedright
  \item \emph{Say it:} \emph{il medioevo}
  \item \emph{Say it:} \emph{il medioevo}, once more
  \item \emph{Say it:} a sentence of your own with \emph{il medioevo}
  \item \emph{Recall:} say the Italian for to free, then the Italian for the age, then say \emph{il medioevo} again
\end{itemize}

\begin{tcolorbox}[breakable,colback=teal!4,colframe=teal!35!black,title={Before you move on}]
What does il medioevo mean? (\textquotedblleft{}The Middle Ages\textquotedblright{}.) Say it once more.
\end{tcolorbox}
\section[le rovine]{le rovine — the ruins}

\label{lesson:IT-C330-rovine}



\begin{quote}
Before the new one: say the Italian for the age, then the Italian for the Middle Ages.
\end{quote}

\subsection*{You'll want to know: le rovine}
\textbf{le rovine} — \textquotedblleft{}the ruins\textquotedblright{}.

\begin{quote}
Abbiamo visitato le rovine di un tempio.
\end{quote}

\begin{quote}
\emph{We visited the ruins of a temple.}
\end{quote}

\emph{Andare in rovina} — to go to ruin.

\begin{cousinweb}
\emph{rovina} is inherited from Latin \emph{ruīna} \textquotedblleft{}a falling down\textquotedblright{}, from \emph{ruere} \textquotedblleft{}to rush, to fall\textquotedblright{} (Lewis \& Short).
\end{cousinweb}

\begin{grammarlens}[title={What you've built}]
One more word for telling, explaining and describing at B1.
\end{grammarlens}

\subsection*{Your turn}
\begin{itemize}
\raggedright
  \item \emph{Say it:} \emph{le rovine}
  \item \emph{Say it:} \emph{le rovine}, once more
  \item \emph{Say it:} a sentence of your own with \emph{le rovine}
  \item \emph{Recall:} say the Italian for the age, then the Italian for the Middle Ages, then say \emph{le rovine} again
\end{itemize}

\begin{tcolorbox}[breakable,colback=teal!4,colframe=teal!35!black,title={Before you move on}]
What does le rovine mean? (\textquotedblleft{}The ruins\textquotedblright{}.) Say it once more.
\end{tcolorbox}
\section[gli scavi]{gli scavi — the excavations, the dig}

\label{lesson:IT-C330-scavi}



\begin{quote}
Before the new one: say the Italian for the Middle Ages, then the Italian for the ruins.
\end{quote}

\subsection*{You'll want to know: gli scavi}
\textbf{gli scavi} — \textquotedblleft{}the excavations, the dig\textquotedblright{}.

\begin{quote}
Negli scavi hanno trovato monete antiche.
\end{quote}

\begin{quote}
\emph{At the dig they found ancient coins.}
\end{quote}

Singular \emph{lo scavo}, built on \emph{scavare}, ‘to dig’.

\begin{cousinweb}
\emph{scavi} is built on \emph{scavare} \textquotedblleft{}to dig out\textquotedblright{}, from Late Latin \emph{excavāre} \textquotedblleft{}to hollow out\textquotedblright{}, built on \emph{cavus} \textquotedblleft{}hollow\textquotedblright{} (Lewis \& Short).
\end{cousinweb}

\begin{grammarlens}[title={What you've built}]
One more word for telling, explaining and describing at B1.
\end{grammarlens}

\subsection*{Your turn}
\begin{itemize}
\raggedright
  \item \emph{Say it:} \emph{gli scavi}
  \item \emph{Say it:} \emph{gli scavi}, once more
  \item \emph{Say it:} a sentence of your own with \emph{gli scavi}
  \item \emph{Recall:} say the Italian for the Middle Ages, then the Italian for the ruins, then say \emph{gli scavi} again
\end{itemize}

\begin{tcolorbox}[breakable,colback=teal!4,colframe=teal!35!black,title={Before you move on}]
What does gli scavi mean? (\textquotedblleft{}The excavations, the dig\textquotedblright{}.) Say it once more.
\end{tcolorbox}
\section[l'archeologo]{l'archeologo — the archaeologist}

\label{lesson:IT-C330-archeologo}



\begin{quote}
Before the new one: say the Italian for the ruins, then the Italian for the excavations.
\end{quote}

\subsection*{You'll want to know: l'archeologo}
\textbf{l'archeologo} — \textquotedblleft{}the archaeologist\textquotedblright{}.

\begin{quote}
L'archeologo ha studiato le rovine per anni.
\end{quote}

\begin{quote}
\emph{The archaeologist studied the ruins for years.}
\end{quote}

Plural \emph{gli archeologi}.

\begin{cousinweb}
\emph{archeologo} is built from Greek \emph{arkhaîos} \textquotedblleft{}ancient\textquotedblright{} and \emph{-lógos} \textquotedblleft{}one who studies\textquotedblright{} (EtymDB has no Italian entry).
\end{cousinweb}

\begin{grammarlens}[title={What you've built}]
One more word for telling, explaining and describing at B1.
\end{grammarlens}

\subsection*{Your turn}
\begin{itemize}
\raggedright
  \item \emph{Say it:} \emph{l'archeologo}
  \item \emph{Say it:} \emph{l'archeologo}, once more
  \item \emph{Say it:} a sentence of your own with \emph{l'archeologo}
  \item \emph{Recall:} say the Italian for the ruins, then the Italian for the excavations, then say \emph{l'archeologo} again
\end{itemize}

\begin{tcolorbox}[breakable,colback=teal!4,colframe=teal!35!black,title={Before you move on}]
What does l'archeologo mean? (\textquotedblleft{}The archaeologist\textquotedblright{}.) Say it once more.
\end{tcolorbox}
